
Podstawą dla przeprowadzenia analizy jest badanie generatora losowości neoTRNG. 
Generatora został zbadany przy użyciu testów losowości: NIST SP 800-22, NIST SP 800-90B 
oraz testu ENT. 

Zbadałem różne aspekty generatora i na podstawie testów próbuję określić 
ich wpływ na jakość genrowanej losowości. Badane aspekty obejmują:

- ilość wykorzystanych oscylatorów pierścieniowych (RO) oraz inwerterów (I) w oscylatorach. 
(Jak skalowanie tych parametrów wpływa na losowość?)

- zastosowanie bądź brak zastosowania korektora von Neumanna i mieszania wielomianem CRC.
(Jak obecność tych metod post-processingu wpływa na losowość?)

- rozmieszczenie RO w układzie FPGA (rozproszone - RO zdala od siebie bądź 
zlokalizowane - RO blisko siebie).
(Jak położenie RO w układzie wpłynie na losowość?)

- zastosowanie innych metod post-processingu: AES-CTR lub ekstraktora Toeplitza.
(Jak te metody post-processing wpłyną na losowość?)

do każdego z pytań muszę wyciągnąć wniosek poparty wynikami w tabelach. Muszę 
stwierdzić czy dany aspekt ma wpływ na losowość, nie ma wpływu na losowość bądź 
wynik nie jest jednoznaczny. 